\documentclass[11pt]{article}
\usepackage[margin=1in]{geometry}
\usepackage{amsmath,amssymb}
\setlength{\parindent}{0pt}
\setlength{\parskip}{6pt}
\begin{document}
\section*{STAT 412 QZ07 --- Question 1}

\subsection*{Problem}

The average height of females in the freshman class of a certain college has historically been 162.2 centimeters with a standard deviation of 6.7 centimeters. Is there reason to believe that there has been a change in the average height if a random sample of 60 females in the present freshman class has an average height of 164.3 centimeters? Use a P-value in your conclusion. Assume the standard deviation remains the same.

\subsection*{Solution}

The claim is that the mean has changed, so this is a two-tailed test:
\[
H_0:\mu=162.2,\qquad H_1:\mu\ne162.2.
\]

Because the population standard deviation is known, use a $z$ test:
\[
z=\frac{\bar{x}-\mu_0}{\sigma/\sqrt{n}}
=\frac{164.3-162.2}{6.7/\sqrt{60}}
=2.43.
\]

The two-tailed P-value is
\[
P=2P(Z\ge 2.43)\approx 0.015.
\]

Since the P-value is small, reject $H_0$ for significance levels greater than 0.015.

\subsection*{Final Answer}

\fbox{\begin{minipage}{0.94\linewidth}
$H_0:\mu=162.2$, $H_1:\mu\ne162.2$; $z=2.43$; P-value $=0.015$. Reject $H_0$. There is sufficient evidence that the average height has changed.
\end{minipage}}

\end{document}
